% assets/a3-az-cases.svg — az_polar 의 네 규칙 중 어느 둘이 요구를 깎는가
%
%   bash _tools/tikz2svg.sh assets/src/a3-az-cases.tex assets/a3-az-cases.svg
%
% 이 그림이 말하는 것
%   (가) 규칙 (b) — 역추력을 금지하면 각도가 쐐기 **밖으로** 밀려난다
%   (나) 규칙 (c) — 각도를 깎으면 크기는 그대로인데 방향이 틀어진다
%   두 경우 모두 오류도 경고도 뜨지 않는다. 깎인 만큼이 못 낸 힘이다
%
% 화면 좌표 — a3-az-polar.tex 와 같은 규칙
%   화면 x = F_y (우현 쪽), 화면 y = F_x (선수 쪽)
%   화면 극각 theta = 90 - delta
%       delta =   0  ->  theta =  90   (위)
%       delta = +45  ->  theta =  45   (오른위)
%       delta = -45  ->  theta = 135   (왼위)
%   쐐기는 theta 가 45~135 인 위쪽 부채꼴
%
% 숫자의 출처 — A3_azpolar_example.m 실행 출력 (한 자리까지 같다)
%   (가) 경우 c:  (F_x, F_y) = (-150, 80),  rev = 0
%        delta_raw = 151.9275  ->  접기  T = -170, delta = -28.0725
%        ->  되돌리기  T = +170, delta = +151.9275  ->  깎기  delta = 45
%        실제 (120.2082, 120.2082),  |e| = 273.1833 N
%   (나) 경우 d:  (F_x, F_y) = (25, 50),  rev = 1
%        T = 55.9017, delta = 63.4349  ->  깎기  delta = 45
%        실제 (39.5285, 39.5285),  |e| = 17.9089 N
%
% 눈금
%   (가) 0.14 mm/N — 170 N 이 23.8 mm
%   (나) 0.40 mm/N — 55.9017 N 이 22.36 mm
%   각도를 깎아도 T 는 그대로이므로 두 화살표의 **길이가 같다**.
%   그래서 반지름이 같은 호 위에서 벡터가 미끄러진다

\documentclass[border=6pt]{standalone}
% capstone-style.tex — 이 볼트의 TikZ 그림이 공유하는 스타일
%
% 쓰는 법:  % capstone-style.tex — 이 볼트의 TikZ 그림이 공유하는 스타일
%
% 쓰는 법:  % capstone-style.tex — 이 볼트의 TikZ 그림이 공유하는 스타일
%
% 쓰는 법:  \input{capstone-style}   (assets/src/ 안에서)
% 빌드:     bash _tools/tikz2svg.sh assets/src/w02-node-topic.tex assets/w02-node-topic.svg
%
% 왜 TikZ 인가
%   손으로 SVG 좌표를 쓰면 화살촉·선 굵기·글자 크기가 그림마다 미묘하게 달라진다.
%   그것이 "자료가 어설퍼 보인다"의 정체다. 스타일을 한 번만 정하고 상대좌표로
%   배치하면 좌표를 고쳐도 그림이 무너지지 않는다.
%   수식은 본문과 같은 TeX 글꼴로 나온다.
%
% 한글
%   ko.TeX(DVI 경로)를 쓴다. XeLaTeX 은 TikZ 그래픽을 PDF special 로 내보내는데
%   dvisvgm 이 그것을 받으려면 Ghostscript < 10.01 이 필요하다. 이 PC 에는
%   10.07 이 깔려 있어 그 경로가 막힌다. latex -> DVI -> dvisvgm 이 유일하게
%   동작하는 조합이다.

\usepackage[hangul]{kotex}
\usepackage{tikz}
\usepackage{amsmath}
\usepackage{xcolor}
\usetikzlibrary{arrows.meta, positioning, calc, fit, backgrounds, decorations.pathmorphing}

% ---- 색 --------------------------------------------------------------------
% 강조는 하나만. 나머지는 검정. 흑백 인쇄에서도 읽혀야 한다.
\definecolor{ink}{HTML}{1A1A1A}
\definecolor{soft}{HTML}{666666}
\definecolor{accent}{HTML}{7C3AED}
\definecolor{warn}{HTML}{B91C1C}

% 계층 색 — Simulink 모델의 블록 색과 같은 규칙을 쓴다
\definecolor{guid}{HTML}{1D4ED8}    % 유도
\definecolor{ctrl}{HTML}{EA580C}    % 제어
\definecolor{thr}{HTML}{D97706}     % 추진기
\definecolor{plant}{HTML}{16A34A}   % 운동모델
\definecolor{nav}{HTML}{7C3AED}     % 항법 · 신호처리
\definecolor{pathc}{HTML}{0D9488}   % 계획 경로 · 항적
\definecolor{errc}{HTML}{C2410C}    % 오차
\definecolor{flow}{HTML}{0369A1}    % 조류 · 바람

% ---- 선과 화살촉 -----------------------------------------------------------
\tikzset{
  sig/.style   = {ink, line width=0.5pt, -{Stealth[length=4pt, width=3pt]}},
  wire/.style  = {ink, line width=0.5pt},
  fb/.style    = {ink, line width=0.5pt, densely dashed,
                  -{Stealth[length=4pt, width=3pt]}},
  asig/.style  = {accent, line width=0.5pt, densely dashed,
                  -{Stealth[length=4pt, width=3pt]}},
  % 블록 — 흰 바탕, 직각, 검은 테두리. 색은 테두리에만 준다
  blk/.style   = {draw=ink, line width=0.55pt, fill=white,
                  minimum width=18mm, minimum height=8.5mm, inner sep=2pt, font=\small},
  gblk/.style  = {blk, draw=guid},
  cblk/.style  = {blk, draw=ctrl},
  tblk/.style  = {blk, draw=thr},
  pblk/.style  = {blk, draw=plant},
  nblk/.style  = {blk, draw=nav},
  % 합산점 · 분기점
  sum/.style   = {draw=ink, line width=0.55pt, fill=white, circle,
                  inner sep=0pt, minimum size=4.6mm, font=\scriptsize},
  dot/.style   = {ink, fill=ink, circle, inner sep=0pt, minimum size=1.5mm},
  % 글자
  sname/.style = {font=\small, inner sep=1.5pt},
  note/.style  = {font=\scriptsize, soft, inner sep=1.5pt},
  wnote/.style = {font=\scriptsize, warn, inner sep=1.5pt},
  anote/.style = {font=\scriptsize, accent, inner sep=1.5pt},
  axis/.style  = {ink, line width=0.5pt},
  tick/.style  = {font=\scriptsize, soft},
  % 묶음 상자 — 서브시스템 표시
  grp/.style   = {draw=soft, line width=0.4pt, densely dotted, rounded corners=1mm,
                  inner sep=3mm},
}
   (assets/src/ 안에서)
% 빌드:     bash _tools/tikz2svg.sh assets/src/w02-node-topic.tex assets/w02-node-topic.svg
%
% 왜 TikZ 인가
%   손으로 SVG 좌표를 쓰면 화살촉·선 굵기·글자 크기가 그림마다 미묘하게 달라진다.
%   그것이 "자료가 어설퍼 보인다"의 정체다. 스타일을 한 번만 정하고 상대좌표로
%   배치하면 좌표를 고쳐도 그림이 무너지지 않는다.
%   수식은 본문과 같은 TeX 글꼴로 나온다.
%
% 한글
%   ko.TeX(DVI 경로)를 쓴다. XeLaTeX 은 TikZ 그래픽을 PDF special 로 내보내는데
%   dvisvgm 이 그것을 받으려면 Ghostscript < 10.01 이 필요하다. 이 PC 에는
%   10.07 이 깔려 있어 그 경로가 막힌다. latex -> DVI -> dvisvgm 이 유일하게
%   동작하는 조합이다.

\usepackage[hangul]{kotex}
\usepackage{tikz}
\usepackage{amsmath}
\usepackage{xcolor}
\usetikzlibrary{arrows.meta, positioning, calc, fit, backgrounds, decorations.pathmorphing}

% ---- 색 --------------------------------------------------------------------
% 강조는 하나만. 나머지는 검정. 흑백 인쇄에서도 읽혀야 한다.
\definecolor{ink}{HTML}{1A1A1A}
\definecolor{soft}{HTML}{666666}
\definecolor{accent}{HTML}{7C3AED}
\definecolor{warn}{HTML}{B91C1C}

% 계층 색 — Simulink 모델의 블록 색과 같은 규칙을 쓴다
\definecolor{guid}{HTML}{1D4ED8}    % 유도
\definecolor{ctrl}{HTML}{EA580C}    % 제어
\definecolor{thr}{HTML}{D97706}     % 추진기
\definecolor{plant}{HTML}{16A34A}   % 운동모델
\definecolor{nav}{HTML}{7C3AED}     % 항법 · 신호처리
\definecolor{pathc}{HTML}{0D9488}   % 계획 경로 · 항적
\definecolor{errc}{HTML}{C2410C}    % 오차
\definecolor{flow}{HTML}{0369A1}    % 조류 · 바람

% ---- 선과 화살촉 -----------------------------------------------------------
\tikzset{
  sig/.style   = {ink, line width=0.5pt, -{Stealth[length=4pt, width=3pt]}},
  wire/.style  = {ink, line width=0.5pt},
  fb/.style    = {ink, line width=0.5pt, densely dashed,
                  -{Stealth[length=4pt, width=3pt]}},
  asig/.style  = {accent, line width=0.5pt, densely dashed,
                  -{Stealth[length=4pt, width=3pt]}},
  % 블록 — 흰 바탕, 직각, 검은 테두리. 색은 테두리에만 준다
  blk/.style   = {draw=ink, line width=0.55pt, fill=white,
                  minimum width=18mm, minimum height=8.5mm, inner sep=2pt, font=\small},
  gblk/.style  = {blk, draw=guid},
  cblk/.style  = {blk, draw=ctrl},
  tblk/.style  = {blk, draw=thr},
  pblk/.style  = {blk, draw=plant},
  nblk/.style  = {blk, draw=nav},
  % 합산점 · 분기점
  sum/.style   = {draw=ink, line width=0.55pt, fill=white, circle,
                  inner sep=0pt, minimum size=4.6mm, font=\scriptsize},
  dot/.style   = {ink, fill=ink, circle, inner sep=0pt, minimum size=1.5mm},
  % 글자
  sname/.style = {font=\small, inner sep=1.5pt},
  note/.style  = {font=\scriptsize, soft, inner sep=1.5pt},
  wnote/.style = {font=\scriptsize, warn, inner sep=1.5pt},
  anote/.style = {font=\scriptsize, accent, inner sep=1.5pt},
  axis/.style  = {ink, line width=0.5pt},
  tick/.style  = {font=\scriptsize, soft},
  % 묶음 상자 — 서브시스템 표시
  grp/.style   = {draw=soft, line width=0.4pt, densely dotted, rounded corners=1mm,
                  inner sep=3mm},
}
   (assets/src/ 안에서)
% 빌드:     bash _tools/tikz2svg.sh assets/src/w02-node-topic.tex assets/w02-node-topic.svg
%
% 왜 TikZ 인가
%   손으로 SVG 좌표를 쓰면 화살촉·선 굵기·글자 크기가 그림마다 미묘하게 달라진다.
%   그것이 "자료가 어설퍼 보인다"의 정체다. 스타일을 한 번만 정하고 상대좌표로
%   배치하면 좌표를 고쳐도 그림이 무너지지 않는다.
%   수식은 본문과 같은 TeX 글꼴로 나온다.
%
% 한글
%   ko.TeX(DVI 경로)를 쓴다. XeLaTeX 은 TikZ 그래픽을 PDF special 로 내보내는데
%   dvisvgm 이 그것을 받으려면 Ghostscript < 10.01 이 필요하다. 이 PC 에는
%   10.07 이 깔려 있어 그 경로가 막힌다. latex -> DVI -> dvisvgm 이 유일하게
%   동작하는 조합이다.

\usepackage[hangul]{kotex}
\usepackage{tikz}
\usepackage{amsmath}
\usepackage{xcolor}
\usetikzlibrary{arrows.meta, positioning, calc, fit, backgrounds, decorations.pathmorphing}

% ---- 색 --------------------------------------------------------------------
% 강조는 하나만. 나머지는 검정. 흑백 인쇄에서도 읽혀야 한다.
\definecolor{ink}{HTML}{1A1A1A}
\definecolor{soft}{HTML}{666666}
\definecolor{accent}{HTML}{7C3AED}
\definecolor{warn}{HTML}{B91C1C}

% 계층 색 — Simulink 모델의 블록 색과 같은 규칙을 쓴다
\definecolor{guid}{HTML}{1D4ED8}    % 유도
\definecolor{ctrl}{HTML}{EA580C}    % 제어
\definecolor{thr}{HTML}{D97706}     % 추진기
\definecolor{plant}{HTML}{16A34A}   % 운동모델
\definecolor{nav}{HTML}{7C3AED}     % 항법 · 신호처리
\definecolor{pathc}{HTML}{0D9488}   % 계획 경로 · 항적
\definecolor{errc}{HTML}{C2410C}    % 오차
\definecolor{flow}{HTML}{0369A1}    % 조류 · 바람

% ---- 선과 화살촉 -----------------------------------------------------------
\tikzset{
  sig/.style   = {ink, line width=0.5pt, -{Stealth[length=4pt, width=3pt]}},
  wire/.style  = {ink, line width=0.5pt},
  fb/.style    = {ink, line width=0.5pt, densely dashed,
                  -{Stealth[length=4pt, width=3pt]}},
  asig/.style  = {accent, line width=0.5pt, densely dashed,
                  -{Stealth[length=4pt, width=3pt]}},
  % 블록 — 흰 바탕, 직각, 검은 테두리. 색은 테두리에만 준다
  blk/.style   = {draw=ink, line width=0.55pt, fill=white,
                  minimum width=18mm, minimum height=8.5mm, inner sep=2pt, font=\small},
  gblk/.style  = {blk, draw=guid},
  cblk/.style  = {blk, draw=ctrl},
  tblk/.style  = {blk, draw=thr},
  pblk/.style  = {blk, draw=plant},
  nblk/.style  = {blk, draw=nav},
  % 합산점 · 분기점
  sum/.style   = {draw=ink, line width=0.55pt, fill=white, circle,
                  inner sep=0pt, minimum size=4.6mm, font=\scriptsize},
  dot/.style   = {ink, fill=ink, circle, inner sep=0pt, minimum size=1.5mm},
  % 글자
  sname/.style = {font=\small, inner sep=1.5pt},
  note/.style  = {font=\scriptsize, soft, inner sep=1.5pt},
  wnote/.style = {font=\scriptsize, warn, inner sep=1.5pt},
  anote/.style = {font=\scriptsize, accent, inner sep=1.5pt},
  axis/.style  = {ink, line width=0.5pt},
  tick/.style  = {font=\scriptsize, soft},
  % 묶음 상자 — 서브시스템 표시
  grp/.style   = {draw=soft, line width=0.4pt, densely dotted, rounded corners=1mm,
                  inner sep=3mm},
}


\begin{document}
\begin{tikzpicture}[x=1mm, y=1mm]

\tikzset{
  axs/.style  = {nav, line width=0.6pt, -{Stealth[length=3.4pt,width=2.4pt]}},
  fvec/.style = {thr, line width=1.4pt, -{Stealth[length=4.4pt,width=3.2pt]}},
  cvec/.style = {warn, line width=1.4pt, -{Stealth[length=4.4pt,width=3.2pt]}},
  evec/.style = {errc, line width=0.9pt, densely dashed,
                 -{Stealth[length=3.6pt,width=2.6pt]}},
  hub/.style  = {circle, draw=thr, fill=white, line width=0.7pt,
                 minimum size=3.4mm, inner sep=0pt},
}

% ===========================================================================
% (가) 규칙 (b) — 역추력 금지가 각도를 쐐기 밖으로 밀어낸다
% ===========================================================================
\begin{scope}[shift={(0,0)}]

  \node[font=\small\bfseries, anchor=north west] at (0,48)
       {(가) \texttt{rev = 0} — 역추력을 금지하면 각도가 쐐기 밖으로 밀려난다};

  \begin{scope}[shift={(44,0)}]

    % --- 앞쪽 쐐기 / 뒤쪽 쐐기 -------------------------------------------
    \fill[guid!10] (0,0) -- (45:30mm) arc (45:135:30mm) -- cycle;
    \fill[errc!8]  (0,0) -- (225:30mm) arc (225:315:30mm) -- cycle;
    \draw[guid, line width=0.4pt] (45:30mm) arc (45:135:30mm);
    \draw[errc, line width=0.4pt, densely dashed] (225:30mm) arc (225:315:30mm);
    \node[font=\scriptsize, guid, anchor=south] at (90:31mm)
         {$T>0$ 로 낼 수 있는 쐐기};
    \node[font=\scriptsize, errc, anchor=north] at (270:31mm)
         {$T<0$ 로 낼 수 있는 쐐기 — \texttt{rev = 0} 이면 금지};

    % --- 좌표축 ----------------------------------------------------------
    \draw[axs] (0,-33) -- (0,34);
    \node[font=\scriptsize, nav, anchor=south] at (0,34.5) {$F_x$ [N]};
    \draw[axs] (-34,0) -- (34,0);
    \node[font=\scriptsize, nav, anchor=west] at (34.5,0) {$F_y$ [N]};

    % --- 요구한 힘 — theta = 90 - 151.9275 = -61.9275, 길이 23.8 --------
    \draw[fvec] (0,0) -- (-61.9275:23.8mm);
    \node[font=\scriptsize, thr, anchor=north west, inner sep=1pt]
         at (-61.9275:24.5mm) {요구 $(-150,\ 80)$};

    % --- 깎인 뒤 실제로 난 힘 — theta = 45, 길이 같음 -------------------
    \draw[cvec] (0,0) -- (45:23.8mm);
    \node[font=\scriptsize, warn, anchor=south west, inner sep=1pt]
         at (45:24.5mm) {실제 $(120.2,\ 120.2)$};

    % --- 같은 반지름의 호 위에서 미끄러진다 ------------------------------
    \draw[warn, line width=0.45pt, densely dotted,
          -{Stealth[length=3pt,width=2.2pt]}]
         (-61.9275:23.8mm) arc (-61.9275:45:23.8mm);

    % --- 못 낸 힘 ---------------------------------------------------------
    \draw[evec] (-61.9275:23.8mm) -- (45:23.8mm);
    \node[font=\scriptsize, errc, anchor=west, inner sep=1pt] at (18,-6)
         {$\lvert \mathbf{e} \rvert = 273.2$ N};

    % --- rev = 1 이면 접은 표기가 쐐기 안이다 ----------------------------
    \draw[plant, line width=1.0pt, densely dashed,
          -{Stealth[length=4pt,width=3pt]}] (0,0) -- (118.0725:21mm);
    \draw[plant, line width=0.35pt] (118.0725:21.5mm) -- (-16,-9);
    \node[font=\scriptsize, plant, anchor=north east, inner sep=1pt, align=right]
         at (-16,-10)
         {\texttt{rev = 1} 이면 여기서 끝 \\ $\delta = -28.07^\circ$ — 쐐기 안};

    \node[hub] at (0,0) {};
  \end{scope}

  \node[anchor=north west, font=\scriptsize, text width=88mm] at (0,-38)
       {\textbf{(a) 접기와 (b) 되돌리기는 같은 힘을 다르게 적을 뿐이다} \\[4pt]
        \hspace*{4mm}접기 전 \quad $T = +170$ N, \ $\delta = +151.93^\circ$ \\[2pt]
        \hspace*{4mm}접은 뒤 \quad $T = -170$ N, \ $\delta = -28.07^\circ$
        \quad $\to$ \ 쐐기 안 \\[2pt]
        \hspace*{4mm}되돌린 뒤 \ $T = +170$ N, \ $\delta = +151.93^\circ$
        \quad $\to$ \ 쐐기 밖 \\[4pt]
        \textcolor{warn}{\textbf{요구를 깎은 것은 (b)가 아니라 그 다음의 (c)다}} \\[2pt]
        되돌리기 자체는 힘을 바꾸지 않는다. 다만 각도를 쐐기 \textbf{밖}
        $151.93^\circ$ 에 놓아 (c)가 $45^\circ$ 로 깎게 만든다 \\[4pt]
        결과: 뒤로 밀어야 할 상황에서 \textbf{앞으로} $120$ N 을 민다.
        부호가 뒤집힌 힘이므로 위치 제어가 발산한다 \\[4pt]
        $\;\to\;$ 그래서 이 과목의 배분기는 \texttt{rev = 1} (역추력 허용)
        을 기본값으로 둔다};
\end{scope}

% ===========================================================================
% (나) 규칙 (c) — 각도를 깎으면 크기는 남고 방향이 틀어진다
% ===========================================================================
\begin{scope}[shift={(98,0)}]

  \node[font=\small\bfseries, anchor=north west] at (0,48)
       {(나) 규칙 (c) — 각도를 깎으면 크기는 그대로, 방향만 틀어진다};

  \begin{scope}[shift={(44,0)}]

    % --- 쐐기 -------------------------------------------------------------
    \fill[guid!10] (0,0) -- (45:30mm) arc (45:135:30mm) -- cycle;
    \draw[guid, line width=0.4pt] (45:30mm) arc (45:135:30mm);
    \draw[soft, line width=0.4pt, densely dashed] (0,0) -- (45:32mm);
    \draw[soft, line width=0.4pt, densely dashed] (0,0) -- (135:32mm);
    \node[font=\scriptsize, guid, anchor=south west] at (45:33mm)
         {$\delta = +45^\circ$ — 한계};
    \node[font=\scriptsize, guid, anchor=south east] at (135:33mm)
         {$\delta = -45^\circ$};

    % --- 좌표축 ----------------------------------------------------------
    \draw[axs] (0,-10) -- (0,34);
    \node[font=\scriptsize, nav, anchor=south] at (0,34.5) {$F_x$ [N]};
    \draw[axs] (-34,0) -- (34,0);
    \node[font=\scriptsize, nav, anchor=west] at (34.5,0) {$F_y$ [N]};

    % --- 요구한 힘 — theta = 90 - 63.4349 = 26.5651, 길이 22.36 ---------
    \draw[fvec] (0,0) -- (26.5651:22.36mm);
    \node[font=\scriptsize, thr, anchor=north west, inner sep=1pt, align=left]
         at (21,10.5) {요구 $(25,\ 50)$ \\ $\delta = 63.43^\circ$};

    % --- 깎인 뒤 — theta = 45, 같은 길이 --------------------------------
    \draw[cvec] (0,0) -- (45:22.36mm);
    \node[font=\scriptsize, warn, anchor=south west, inner sep=1pt, align=left]
         at (45:23.5mm) {실제 $(39.53,\ 39.53)$ \\ $\delta = 45^\circ$ — 한계에 붙음};

    % --- 같은 반지름 호 ---------------------------------------------------
    \draw[warn, line width=0.45pt, densely dotted,
          -{Stealth[length=3pt,width=2.2pt]}]
         (26.5651:22.36mm) arc (26.5651:45:22.36mm);
    \node[font=\scriptsize, warn, anchor=north west, inner sep=1pt, align=left]
         at (24,-4) {$T = 55.90$ N 고정 \\ — 화살표 길이가 같다};

    % --- 못 낸 힘 ---------------------------------------------------------
    \draw[evec] (26.5651:22.36mm) -- (45:22.36mm);
    \draw[errc, line width=0.35pt] (18.5,13) -- (9,13);
    \node[font=\scriptsize, errc, anchor=east, inner sep=1pt] at (8.5,13)
         {$\lvert \mathbf{e} \rvert = 17.91$ N};

    \node[hub] at (0,0) {};
  \end{scope}

  \node[anchor=north west, font=\scriptsize, text width=88mm] at (0,-38)
       {\textbf{깎기는 벡터를 호 위에서 미끄러뜨린다} \\[4pt]
        $T$ 는 건드리지 않으므로 화살표 \textbf{길이가 같다}.
        방향만 한계선에 붙는다 \\[4pt]
        \hspace*{4mm}요구 \quad $F_x = 25$, \ $F_y = 50$ \\[2pt]
        \hspace*{4mm}실제 \quad $F_x = 39.53$, \ $F_y = 39.53$ \\[2pt]
        \hspace*{4mm}차이 \quad $F_x$ 는 \textbf{더} 나고, $F_y$ 는 \textbf{덜} 난다 \\[4pt]
        \textcolor{warn}{\textbf{요구하지 않은 축까지 늘어나므로 유의}} \\[2pt]
        전진력을 요구하지 않았는데 $14.53$ N 이 앞으로 더 났다.
        옆으로 가려는 배가 앞으로 밀린다 \\[4pt]
        $\;\to\;$ 이 $(25,\ 50)$ 이 1-7절 예제의 \textbf{우현 추진기}다.
        null space 보정은 깎기가 일어나기 \textbf{전에}
        $(25,\ 25)$ 로 바꿔 놓아 이 손실을 없앤다};
\end{scope}

\end{tikzpicture}
\end{document}
